% =====================================================
% 题型：四边形分割与海伦公式综合解答题 (q_triangle_heron_quadrilateral.tex)
% =====================================================
% 难度：★★★ (难题)
% =====================================================
\newcommand{\qTriangleHeronQuadrilateralStem}{在四边形 \(ABCD\) 中，\(AB=BC=CD=2\)，\(AD=4\)，\(\angle ABC=120^\circ\)，求四边形 \(ABCD\) 面积的最大值。}

\newcommand{\qTriangleHeronQuadrilateralFigure}[1][0.9]{%
  \begin{tikzpicture}[scale=#1, line join=round, line cap=round]
    \coordinate (A) at (0,0);
    \coordinate (B) at (2,0);
    \coordinate (C) at (3,1.732);
    \coordinate (D) at (0.3,2.2);
    \draw[thick] (A)--(B)--(C)--(D)--cycle;
    \draw[dashed] (A)--(C);
    \node[below left] at (A) {$A$};
    \node[below right] at (B) {$B$};
    \node[right] at (C) {$C$};
    \node[above left] at (D) {$D$};
  \end{tikzpicture}%
}

\newcommand{\qTriangleHeronQuadrilateralAnswer}{\(3\sqrt3+2\sqrt7\)}

\newcommand{\qTriangleHeronQuadrilateralSolution}{%
  连接对角线 \(AC\)。在 \(\triangle ABC\) 中，由余弦定理：
  \[
    AC^2=AB^2+BC^2-2AB\cdot BC\cos120^\circ
    =4+4-8\left(-\frac12\right)=12.
  \]
  所以
  \[
    AC=2\sqrt3.
  \]
  同时
  \[
    S_{\triangle ABC}
    =\frac12\cdot2\cdot2\cdot\sin120^\circ
    =\sqrt3.
  \]

  四边形面积为
  \[
    S_{ABCD}=S_{\triangle ABC}+S_{\triangle ACD}.
  \]
  其中 \(\triangle ACD\) 的三边为
  \[
    AC=2\sqrt3,
    \qquad CD=2,
    \qquad AD=4.
  \]
  这三边长度均固定，因此 \(\triangle ACD\) 实际上由 SSS 唯一确定，面积可直接用海伦公式计算。

  半周长：
  \[
    p=\frac{2\sqrt3+2+4}{2}=3+\sqrt3.
  \]
  由海伦公式：
  \[
    S_{\triangle ACD}
    =\sqrt{p(p-2\sqrt3)(p-2)(p-4)}.
  \]
  代入并化简可得
  \[
    S_{\triangle ACD}=2\sqrt7.
  \]
  因此
  \[
    S_{ABCD}=\sqrt3+2\sqrt7.
  \]

  若题目原始设问确为“面积最大值”，则还需要核对四边形顶点顺序与可翻折构型是否允许另一面积分支；按当前凸四边形构型，面积由边长数据唯一确定。%
}

\newcommand{\qTriangleHeronQuadrilateralDiagnosis}{%
  高频问题是看到“四边形面积最大值”就直接设角求导，却没有先连接对角线检查两块三角形是否已经被固定。应先做自由度判断，再决定是否真的存在连续最值变量。%
}

\newcommand{\qTriangleHeronQuadrilateralTeachingNotes}{%
  本题适合训练“先分割、再数自由度”。连接 \(AC\) 后，\(\triangle ABC\) 可由 SAS 固定，进而 \(AC\) 固定；随后 \(\triangle ACD\) 又由 SSS 固定。若原卷答案声称还有连续最值，应优先回查原始题面。%
}
